
\section{量子力学：Pöschl--Teller 势的低能散射——精确解与相移提取}
\label{sec:poschl-teller}

\begin{modelbox}[Pöschl--Teller 势的 s 波散射]
质量为 $m$ 的无自旋粒子，受到中心势场
\[
V(r) = -\frac{\hbar^2}{ma^2}\frac{1}{\cosh^2(r/a)},\qquad a>0
\]
的作用。

已知方程
\[
\frac{d^2y}{dx^2} + K^2 y + \frac{2}{\cosh^2 x}\,y = 0
\]
的两个线性独立解为
\[
y_1(x) = e^{iKx}(\tanh x - iK),\qquad
y_2(x) = e^{-iKx}(\tanh x + iK).
\]

在低能极限下，求粒子能量为 $E$ 时 s 波的散射截面及其角分布。
\end{modelbox}

\subsection{题型反射弧标签}

\begin{tipbox}[看到什么→启动什么]
\begin{center}
\renewcommand{\arraystretch}{1.5}
\begin{tabular}{|c|c|l|}
\hline
\textbf{一级（关键词）} & \textbf{二级（方法）} & \textbf{三级（操作）} \\
\hline
"已知精确解"、"已给两个解" & 直接代入边界条件 & 不需要从头解方程 \\
\hline
$V(r) \propto 1/\cosh^2(r/a)$ & Pöschl--Teller 势 & 记住形式 $V = -V_0/\cosh^2(r/a)$ \\
\hline
低能、$l=0$、s 波 & 只保留 s 波 & $\sigma \approx 4\pi/k^2 \sin^2\delta_0$ \\
\hline
"角分布"、"微分截面" & s 波各向同性 & $d\sigma/d\Omega = \sigma_0/4\pi$ \\
\hline
\end{tabular}
\end{center}
\end{tipbox}

\subsection{标准解题过程}

\begin{derivationbox}[径向方程与无量纲化]
\textbf{Step 1：写出径向方程（$l=0$）}

令 $\psi(r) = u(r)/r$，则 $u(r)$ 满足
\[
-\frac{\hbar^2}{2m}\frac{d^2u}{dr^2} + V(r)u = Eu,
\]
即
\[
\frac{d^2u}{dr^2} + \frac{2m}{\hbar^2}\left[E + \frac{\hbar^2}{ma^2}\frac{1}{\cosh^2(r/a)}\right]u = 0.
\]

\textbf{Step 2：无量纲化}

令
\[
x = \frac{r}{a},\quad u(r) = y(x),\quad K = ka,\quad k = \sqrt{\frac{2mE}{\hbar^2}}.
\]

代入得
\[
\frac{d^2y}{dx^2} + K^2 y + \frac{2}{\cosh^2 x}\,y = 0.
\]
这正是题目给出的标准形式。
\end{derivationbox}

\begin{derivationbox}[边界条件与系数确定]
\textbf{Step 3：写出通解}

\[
y(x) = A\,e^{iKx}(\tanh x - iK) + B\,e^{-iKx}(\tanh x + iK).
\]

回到原变量：
\[
u(r) = A\,e^{ikr}\bigl[\tanh(r/a) - ika\bigr] + B\,e^{-ikr}\bigl[\tanh(r/a) + ika\bigr].
\]

\textbf{Step 4：原点边界条件 $u(0)=0$}

对 s 波，$\psi(r) = u(r)/r$ 在原点有限要求 $u(0)=0$。

在 $r=0$ 处，$\tanh(0)=0$，因此
\[
u(0) = A(-ika) + B(ika) = ika(-A+B) = 0\quad\Rightarrow\quad A = B.
\]

于是
\[
u(r) = A\Bigl\{e^{ikr}\bigl[\tanh(r/a) - ika\bigr] + e^{-ikr}\bigl[\tanh(r/a) + ika\bigr]\Bigr\}.
\]
\end{derivationbox}

\begin{derivationbox}[渐近行为与相移提取]
\textbf{Step 5：$r\to\infty$ 的渐近形式}

当 $r\to\infty$ 时，$\tanh(r/a)\to 1$，因此
\begin{align*}
u(r) &\to A\Bigl\{e^{ikr}(1-ika) + e^{-ikr}(1+ika)\Bigr\}\\
&= A\Bigl\{(e^{ikr}+e^{-ikr}) - ika(e^{ikr}-e^{-ikr})\Bigr\}\\
&= 2A\bigl(\cos kr + ka\sin kr\bigr).
\end{align*}

\textbf{Step 6：与标准散射渐近式对比}

标准形式：
\[
u(r) \xrightarrow{r\to\infty} \sin(kr+\delta_0) = \sin\delta_0\cos kr + \cos\delta_0\sin kr.
\]

对比 $\cos kr$ 与 $\sin kr$ 的系数：
\[
\cos kr + ka\sin kr = C\bigl(\sin\delta_0\cos kr + \cos\delta_0\sin kr\bigr).
\]

得到
\[
C\sin\delta_0 = 1,\qquad C\cos\delta_0 = ka.
\]

两式相除：
\[
\boxed{\cot\delta_0 = ka = \sqrt{\frac{2mE}{\hbar^2}}\,a}.
\]
\end{derivationbox}

\begin{formulabox}[散射截面与角分布]
由 $\cot\delta_0 = ka$，得
\[
\sin^2\delta_0 = \frac{1}{1+\cot^2\delta_0} = \frac{1}{1+k^2a^2}.
\]

s 波散射截面：
\[
\sigma_0 = \frac{4\pi}{k^2}\sin^2\delta_0 = \frac{4\pi}{k^2}\cdot\frac{1}{1+k^2a^2}
= \frac{4\pi}{k^2(1+k^2a^2)}.
\]

用 $k^2 = 2mE/\hbar^2$ 代回：
\[
\boxed{\sigma_0 = \frac{2\pi\hbar^2}{mE}\cdot\frac{1}{1+\dfrac{2mEa^2}{\hbar^2}}}.
\]

\textbf{低能极限}（$k\to 0$）：s 波占绝对主导，高阶分波可忽略。

\textbf{角分布}：s 波散射各向同性，
\[
\boxed{\frac{d\sigma}{d\Omega} = \frac{\sigma_0}{4\pi}}.
\]
\end{formulabox}

\subsection{物理图像与补充说明}

\begin{insightbox}[为什么 Pöschl--Teller 势有漂亮解析解？]
\textbf{本质原因}：这个势是量子力学中少数几个\textbf{形状不变}（shape-invariant）势之一，与超对称量子力学（SUSY QM）和可积系统密切相关。

\textbf{关键观察}：方程
\[
\frac{d^2y}{dx^2} + \left[K^2 + \frac{2}{\cosh^2 x}\right]y = 0
\]
可以通过变量替换 $z = \tanh x$ 化为\textbf{连带 Legendre 方程}的形式，其解自然包含 $\tanh x$ 和指数因子。

\textbf{更一般的形式}：Pöschl--Teller 势的标准形式为
\[
V(x) = -\frac{\hbar^2\alpha^2}{2m}\frac{\lambda(\lambda-1)}{\cosh^2\alpha x},
\]
当 $\lambda=2$、$\alpha=1/a$ 时，即为本题的势。该势的反射系数和透射系数都可以精确计算，是散射理论中最重要的\textbf{可精确求解模型}之一。
\end{insightbox}

\begin{tipbox}[考场速判：低能散射的通用套路]
\begin{enumerate}[nosep]
\item \textbf{确认 $l=0$ 主导}：低能时 $k\to 0$，离心势垒 $l(l+1)\hbar^2/(2mr^2)$ 压制高阶分波，只需考虑 s 波。
\item \textbf{写出标准渐近式}：$u(r)\to \sin(kr+\delta_0)$。
\item \textbf{提取相移}：将精确解（或近似解）的渐近形式与标准式对比，读出 $\delta_0$。
\item \textbf{计算截面}：$\sigma_0 = 4\pi/k^2 \sin^2\delta_0$。
\item \textbf{判断角分布}：s 波 $\Rightarrow$ 各向同性；p 波 $\Rightarrow$ $\cos^2\theta$；更高阶 $\Rightarrow$ Legendre 多项式叠加。
\end{enumerate}
\end{tipbox}

\begin{warnbox}[常见错误与思路短路]
\textbf{错误 1}：忘记 $u(0)=0$ 的边界条件，直接令 $A,B$ 为任意常数。

\textbf{错误 2}：渐近行为提取时，把 $\tanh(r/a)$ 的极限搞错——$r\to\infty$ 时 $\tanh\to\mathbf{1}$，不是 0。

\textbf{错误 3}：相移提取时，没有将 $2A(\cos kr + ka\sin kr)$ 写成 $\sin(kr+\delta_0)$ 的标准形式，而是用 $\cos(kr+\delta_0)$ 对比，导致符号混乱。

\textbf{错误 4}：低能极限下忘记只保留 s 波，误以为还需要计算 p 波、 d 波的贡献。

\textbf{记忆锚点}：
\[
\cot\delta_0 = ka\quad\Longleftrightarrow\quad \sigma_0 = \frac{4\pi}{k^2}\cdot\frac{1}{1+k^2a^2}.
\]
对于吸引性的 $1/\cosh^2$ 势，相移 $\delta_0 > 0$（因为势能把粒子"拉进来"，等效于波函数向内移动）。
\end{warnbox}
